\documentclass[10pt,a4paper]{amsart}
\usepackage[margin=19mm]{geometry}
\usepackage{amsmath,amssymb,mathtools}
\usepackage[scr=rsfs,cal=euler]{mathalfa}
\usepackage{xfrac}
\usepackage[unicode,hidelinks,bookmarks=false]{hyperref}
\newcommand{\ie}{i.e.\ }
\newcommand{\cf}{cf.\ }
\newcommand{\resp}{respectively\ }
\newcommand{\Subsection}{\subsection}
\providecommand{\nobreakdash}{\nobreak\hbox{-}}
\setlength{\emergencystretch}{2em}
\multlinegap0pt
\usepackage{bm}
\usepackage{bbm}
\usepackage{dsfont}
\usepackage{xspace}
\usepackage{cancel}
\usepackage{mathtools}
\usepackage{amsthm}
\makeatletter
\@ifundefined{corollary}{\newtheorem{corollary}{Corollary}}{}
\@ifundefined{theorem}{\newtheorem{theorem}{Theorem}}{}
\makeatother
\title[A Dilogarithm Series and a Generalization of a Theorem of Br]{A Dilogarithm Series and a Generalization of a Theorem of Bridgeman}
\author{Chance Sanford}
\date{Source: arXiv:2410.21619v2 (CC BY 4.0)}
\begin{document}
\maketitle

\noindent\emph{Source and licence.} A Dilogarithm Series and a Generalization of a Theorem of Bridgeman. Version arXiv:2410.21619v2 (CC BY 4.0), distributed under the Creative Commons Attribution 4.0 International licence, as recorded in the arXiv licence metadata for this exact version and shown on the abstract page https://arxiv.org/abs/2410.21619v2. Full rights notice: licenses/arxiv-2410.21619-CC-BY-4.0.txt.\par\smallskip

\noindent\textit{Editorial scope.} Counted unit: the Theorem whose source label is \texttt{thm:generalized\_bridgeman} in the article (source0.tex lines 302--311, source label \texttt{thm:generalized\_bridgeman}), delivered together with the article's own complete original proof of it (source0.tex lines 312--389, one \texttt{proof} environment of 78 lines). No mathematics is written, repaired or re-derived: statement and proof are the article's own text line for line. Required context retained uncounted: cor:main (lines 278--283). Every definition, display and label of those lines is retained unchanged. Omitted from this package and not redistributed: 1--277, 284--301, 390--597. Nothing inside a retained excerpt is omitted or altered: every retained body is delivered exactly as the article's lines are written, so no omission receipt of any kind accompanies this package. Citations: the retained excerpts carry no citation command at all, so no bibliography entry of the article is transcribed and no reference is redistributed. Reference adaptation: none was needed and none was made; every reference inside the delivered unit resolves inside it, so no unresolved reference or citation is produced. Printed slips kept literal: no printed word, symbol, hypothesis or number of the counted statement or of its proof was altered. Layout and class: the article's own class is \texttt{article} with packages including amssymb, latexsym, amsmath, epsfig, amsthm, mathtools; the wrapper is the lane's pinned \texttt{amsart} editorial wrapper at 10pt on A4, which restates the theorem-like environments the retained text needs and adds the article's own macro definitions that the retained text needs (none), with extra wrapper packages (bm, bbm, dsfont, xspace, cancel, mathtools). The delivered numbering is the unit's own. Assets: no retained span contains a figure environment, a graphics-inclusion command, a PDF-page inclusion, a table or a TikZ picture; no image file is copied into this package, the package declares no asset dependency and no image file is redistributed with it. Licence: version arXiv:2410.21619v2 of this article is distributed under the Creative Commons Attribution 4.0 International licence, as recorded in the arXiv licence metadata for this exact version and shown on the abstract page https://arxiv.org/abs/2410.21619v2. A full rights notice is delivered at licenses/arxiv-2410.21619-CC-BY-4.0.txt. Characters: every retained excerpt is pure ASCII, so the delivered engine, XeLaTeX with its default Latin Modern text font, reports no missing character.
\begin{corollary}\label{cor:main}
    If $a \in (0,1)$, then 
    \begin{equation*}
        \sum_{n=1}^{\infty}\mathcal{L}\left(\frac{a^2\left(\frac{1}{2}-\frac{a}{2}\right)^{n}\left(\frac{1}{2} + \frac{a}{2}\right)^{n}}{\left(\left(\frac{1}{2}+\frac{a}{2}\right)^{n+1}-\left(\frac{1}{2} - \frac{a}{2}\right)^{n+1}\right)^2}\right) = \mathcal{L}\left(\frac{1-a}{1+a}\right).
    \end{equation*}  
\end{corollary}

\begin{theorem}\label{thm:generalized_bridgeman}
    If $Q > 0$ and $k \in \mathbb{Z}_{>0}$, then
    \begin{equation}
        \sum_{n=1}^{\infty}\mathcal{L}\left(\frac{U^2_k Q^{kn}}{U^2_{k(n+1)}}\right) = \mathcal{L}\left(\frac{Q^k}{\alpha^{2k}}\right).\label{eq:lucas_series_pos}
    \end{equation}
    If $Q < 0$ and $k \in \mathbb{Z}_{>0}$ is an odd integer, then 
    \begin{equation}
       \sum_{n=1}^{\infty}\mathcal{L}\left(-\frac{V^2_k Q^{k(2n-1)}}{D U^2_{2kn} }\right) + \sum_{n=1}^{\infty}\mathcal{L}\left(\frac{V^2_k Q^{2kn}}{ V^2_{k(2n+1)}}\right) = \mathcal{L}\left(-\frac{Q^k}{\alpha^{2k}}\right).\label{eq:lucas_series_neg}
    \end{equation} 
\end{theorem}

\begin{proof}
To prove Equation~\eqref{eq:lucas_series_pos}, notice that since $Q > 0$ and $D = P^2 - 4Q > 0$, $\beta = (P-\sqrt{D})/2$ is also positive. 
 Therefore, the quotient
\begin{equation*}
    \frac{\sqrt{D} U_k}{V_k} = \frac{\alpha^k - \beta^k}{\alpha^k + \beta^k}
\end{equation*}
    is less than 1 for all $k \in \mathbb{Z}_{> 0}$.  Thus, we are permitted to set $a = U_k\sqrt{D}/V_k$ in Corollary \ref{cor:main}.  In doing so, we find that
    \begin{equation}
        \sum_{n=1}^{\infty}\mathcal{L}\left(\frac{U^2_k D\left(\frac{1}{2}-\frac{U_k\sqrt{D}}{2V_k}\right)^{n}\left(\frac{1}{2} + \frac{U_k\sqrt{D}}{2 V_k}\right)^{n}}{V^2_k\left(\left(\frac{1}{2}+\frac{U_k\sqrt{D}}{2 V_k}\right)^{n+1}-\left(\frac{1}{2} - \frac{U_k \sqrt{D}}{2 V_k}\right)^{n+1}\right)^2}\right) = \mathcal{L}\left(\frac{V_k-U_k \sqrt{D}}{V_k+U_k\sqrt{D}}\right).\label{eq:lucas_proof_1}
    \end{equation}
    Examining the left side of Equation \eqref{eq:lucas_proof_1} first, we see that the numerator of the dilogarithm's argument may be expressed as
    \begin{align*}
        U^2_k D\left(\tfrac{1}{2}-\tfrac{U_k\sqrt{D}}{2V_k}\right)^{n}\left(\tfrac{1}{2} + \tfrac{U_k\sqrt{D}}{2 V_k}\right)^{n} &= U^2_k D V^{-2n}_k\left(\tfrac{V_k-U_k\sqrt{D}}{2}\right)^{n}\left(\tfrac{V_k+U_k\sqrt{D}}{2}\right)^{n} \\
        &= U^2_k D V^{-2n}_k \beta^{kn} \alpha^{kn} \\
        &= U^2_k D V^{-2n}_k Q^{kn},
    \end{align*}
    while the denominator may be written as
    \begin{align*}
        \MoveEqLeft[12]  V^2_k\left(\left(\tfrac{1}{2}+\tfrac{U_k\sqrt{D}}{2 V_k}\right)^{n+1} -\left(\tfrac{1}{2} - \tfrac{U_k \sqrt{D}}{2  V_k}\right)^{n+1}\right)^2 \\ 
        ={}& V^{-2n}_k\left(\left(\tfrac{V_k + U_k\sqrt{D}}{2}\right)^{n+1}-\left(\tfrac{V_k - U_k \sqrt{D}}{2}\right)^{n+1}\right)^2 \\
        ={}& V^{-2n}_k \left(\alpha^{k(n+1)} - \beta^{k(n+1)}\right)^2.
    \end{align*}
    Therefore, the sum simplifies to
    \begin{equation*}
        \sum_{n=1}^{\infty}\mathcal{L}\left(\frac{U^2_k D Q^{kn}}{(\alpha^{k(n+1)} - \beta^{k(n+1)})^2}\right) =  \sum_{n=1}^{\infty}\mathcal{L}\left(\frac{U^2_k Q^{kn}}{U^2_{k(n+1)}}\right).
    \end{equation*}
    
    Turning our attention to the right side of Equation~\eqref{eq:lucas_proof_1}, we see that 
    \begin{align*}
        \mathcal{L}\left(\frac{V_k-U_k \sqrt{D}}{V_k+U_k\sqrt{D}}\right) &= \mathcal{L}\left(\frac{V^2_k-U^2_k D}{(V_k+U_k\sqrt{D})^2}\right) \\
        &= \mathcal{L}\left(\frac{4Q^k}{(V_k+U_k\sqrt{D})^2}\right) \\
        &= \mathcal{L}\left(\frac{Q^k}{\alpha^{2k}}\right).
    \end{align*}
    Thus, we have shown that
    \begin{equation*}
        \sum_{n=1}^{\infty}\mathcal{L}\left(\frac{U^2_k Q^{kn}}{U^2_{k(n+1)}}\right)  = \mathcal{L}\left(\frac{Q^k}{\alpha^{2k}}\right),
    \end{equation*}
   proving Equation~\eqref{eq:lucas_series_pos}.
   
   Moving on to Equation~\eqref{eq:lucas_series_neg}, since $Q < 0$ and $k$ is an odd integer, we see that $\beta^k < 0$.  Therefore, the quotient
\begin{equation*}
    \frac{V_k}{\sqrt{D} U_k} = \frac{\alpha^k + \beta^k}{\alpha^k - \beta^k}
\end{equation*}
is less than 1.  Setting $a = V_k/U_k \sqrt{D}$ in Corollary \ref{cor:main}, we find that
    \begin{equation*}
        \sum_{n=1}^{\infty}\mathcal{L}\left(\frac{V^2_k\left(\frac{1}{2}-\frac{V_k}{2U_k\sqrt{D}}\right)^{n}\left(\frac{1}{2} + \frac{V_k}{2U_k\sqrt{D}}\right)^{n}}{U^2_k D\left(\left(\frac{1}{2}+\frac{V_k}{2U_k\sqrt{D}}\right)^{n+1}-\left(\frac{1}{2} - \frac{V_k}{2U_k\sqrt{D}}\right)^{n+1}\right)^2}\right) = \mathcal{L}\left(\frac{U_k\sqrt{D}-V_k}{U_k\sqrt{D}+V_k}\right).
    \end{equation*}
    Once again, simplifying the numerator and denominator separately, the numerator is seen to be
    \begin{equation*}
        V^2_k\left(\tfrac{1}{2}-\tfrac{V_k}{2U_k\sqrt{D}}\right)^{n}\left(\tfrac{1}{2} + \tfrac{V_k}{2U_k\sqrt{D}}\right)^{n} = (-1)^n U^{-2n}_k D^{-n} V^2_k Q^{kn},
    \end{equation*}
    while the denominator admits the form
    \begin{multline*}
     U^2_k D\left(\left(\tfrac{1}{2}+\tfrac{V_k}{2U_k\sqrt{D}}\right)^{n+1}-\left(\tfrac{1}{2} - \tfrac{V_k}{2U_k\sqrt{D}}\right)^{n+1}\right)^2 \quad\quad \quad\quad \quad\quad \quad\quad {}\\
     \begin{aligned}
        &= U^{-2n}D^{-n}\left(\left(\tfrac{V_k+U_k\sqrt{D}}{2}\right)^{n+1}-(-1)^{n+1}\left(\tfrac{V_k-U_k\sqrt{D}}{2} \right)^{n+1}\right)^2 \\
        &= U^{-2n} D^{-n} (\alpha^{k(n+1)} + (-1)^{n} \beta^{k(n+1)})^2.
     \end{aligned}
    \end{multline*}
    Putting this all together yields 
    \begin{equation*}
        \sum_{n=1}^{\infty}\mathcal{L}\left(\frac{V^2_k (-1)^n Q^{kn}}{\left(\alpha^{k(n+1)}+(-1)^{n}\beta^{k(n+1)}\right)^2 }\right) = \mathcal{L}\left(\frac{U_k\sqrt{D}-V_k}{U_k\sqrt{D}+V_k}\right).
    \end{equation*}
    
    Next, by splitting the series into two sums, one over all even indices and the other over all odd indices, we find that
    \begin{multline*}
        \sum_{n=1}^{\infty}\mathcal{L}\left(\frac{V^2_k (-1)^n Q^{kn}}{\left(\alpha^{k(n+1)}+(-1)^{n}\beta^{k(n+1)}\right)^2 }\right) \\ = \sum_{n=1}^{\infty}\mathcal{L}\left(-\frac{V^2_k Q^{k(2n-1)}}{\left(\alpha^{2kn}-\beta^{2kn}\right)^2 }\right) + \sum_{n=1}^{\infty}\mathcal{L}\left(\frac{V^2_k Q^{2kn}}{\left(\alpha^{k(2n+1)}+\beta^{k(2n+1)}\right)^2 }\right).
    \end{multline*}
    Using the Binet formulas for $U_n$ and $V_n$, we arrive at
        \begin{multline*}
        \sum_{n=1}^{\infty}\mathcal{L}\left(\frac{V^2_k (-1)^n Q^{kn}}{\left(\alpha^{k(n+1)}+(-1)^{n}\beta^{k(n+1)}\right)^2 }\right) \\ = \sum_{n=1}^{\infty}\mathcal{L}\left(-\frac{V^2_k Q^{k(2n-1)}}{D U^2_{2kn} }\right) + \sum_{n=1}^{\infty}\mathcal{L}\left(\frac{V^2_k Q^{2kn}}{ V^2_{k(2n+1)}}\right).
    \end{multline*}
    Finally, the right-hand side may be simplified in the same way we proceeded in the proof of Equation~\eqref{eq:lucas_series_pos} to show that
    \begin{equation*}
        \mathcal{L}\left(\frac{U_k\sqrt{D}-V_k}{U_k\sqrt{D}+V_k}\right) = \mathcal{L}\left(-\frac{Q^k}{\alpha^{2k}}\right).
    \end{equation*}
    This completes the proof of Theorem~\ref{thm:generalized_bridgeman}.
\end{proof}

\end{document}
